\section{Functoriality of categories of presheaves}
\label{section-functoriality-PSh}

\noindent
Let $u : \mathcal{C} \to \mathcal{D}$ be a functor between categories.
In this case we denote
$$
u^p :
\textit{PSh}(\mathcal{D})
\longrightarrow
\textit{PSh}(\mathcal{C})
$$
the functor that associates to $\mathcal{G}$ on $\mathcal{D}$ the presheaf
$u^p\mathcal{G} = \mathcal{G} \circ u$. Note that by the previous section
this functor commutes with all limits and colimits.

\medskip\noindent
For $V \in \Ob(\mathcal{D})$ let $\mathcal{I}^u_V$
denote the category with
\begin{equation}
\label{equation-colim-category}
\begin{matrix}
\Ob(\mathcal{I}^u_V)
&
=
&
\{
(U, \phi)
\mid
U \in \Ob(\mathcal{C}),
\phi : V \to u(U)
\}
\\
\Mor_{\mathcal{I}^u_V}((U, \phi), (U', \phi'))
&
=
&
\{
f : U \to U' \text{ in }\mathcal{C}
\mid
u(f) \circ \phi = \phi'
\}
\end{matrix}
\end{equation}
We sometimes drop the superscript ${}^u$ from the notation and we simply write
$\mathcal{I}_V$.
We will use these categories to define a left adjoint to the functor $u^p$.
Before we do so we prove a few technical lemmas.

\begin{lemma}
\label{lemma-almost-directed}
Let $u : \mathcal{C} \to \mathcal{D}$ be a functor between categories.
Suppose that $\mathcal{C}$ has fibre products and equalizers, and that
$u$ commutes with them. Then the categories $(\mathcal{I}_V)^{opp}$
satisfy the hypotheses of
Categories, Lemma \ref{categories-lemma-split-into-directed}.
\end{lemma}

\begin{proof}
There are two conditions to check.

\medskip\noindent
First, suppose we are given three objects
$\phi : V \to u(U)$, $\phi' : V \to u(U')$, and $\phi'' : V \to u(U'')$
and morphisms $a : U' \to U$, $b : U'' \to U$ such that
$u(a) \circ \phi' = \phi$ and $u(b) \circ \phi'' = \phi$.
We have to show there exists another object $\phi''' : V \to u(U''')$
and morphisms $c : U''' \to U'$ and $d : U''' \to U''$ such that
$u(c) \circ \phi''' = \phi'$, $u(d) \circ \phi''' = \phi''$ and
$a \circ c = b \circ d$. We take $U''' = U' \times_U U''$
with $c$ and $d$ the projection morphisms. This works as $u$ commutes
with fibre products; we omit the verification.

\medskip\noindent
Second, suppose we are given two objects
$\phi : V \to u(U)$ and $\phi' : V \to u(U')$
and morphisms $a, b : (U, \phi) \to (U', \phi')$.
We have to find a morphism $c : (U'', \phi'') \to (U, \phi)$
which equalizes $a$ and $b$. Let $c : U'' \to U$ be the equalizer of
$a$ and $b$ in the category $\mathcal{C}$. As $u$ commutes
with equalizers and since $u(a) \circ \phi = u(b) \circ \phi = \phi'$
we obtain a morphism $\phi'' : V \to u(U'')$.
\end{proof}

\begin{lemma}
\label{lemma-directed}
Let $u : \mathcal{C} \to \mathcal{D}$ be a functor between categories.
Assume
\begin{enumerate}
\item the category $\mathcal{C}$ has a final object $X$ and
$u(X)$ is a final object of $\mathcal{D}$ , and
\item the category $\mathcal{C}$ has fibre products and
$u$ commutes with them.
\end{enumerate}
Then the index categories $(\mathcal{I}^u_V)^{opp}$ are filtered (see
Categories, Definition \ref{categories-definition-directed}).
\end{lemma}

\begin{proof}
The assumptions imply that the assumptions of
Lemma \ref{lemma-almost-directed}
are satisfied (see the discussion in
Categories, Section \ref{categories-section-finite-limits}).
By
Categories, Lemma \ref{categories-lemma-split-into-directed}
we see that $\mathcal{I}_V$ is a (possibly empty) disjoint union of
directed categories.
Hence it suffices to show that $\mathcal{I}_V$ is connected.

\medskip\noindent
First, we show that $\mathcal{I}_V$ is nonempty.
Namely, let $X$ be the final object of $\mathcal{C}$,
which exists by assumption.
Let $V \to u(X)$ be the morphism coming from the fact
that $u(X)$ is final in $\mathcal{D}$ by assumption.
This gives an object of $\mathcal{I}_V$.

\medskip\noindent
Second, we show that $\mathcal{I}_V$ is connected.
Let $\phi_1 : V \to u(U_1)$ and $\phi_2 : V \to u(U_2)$ be
in $\Ob(\mathcal{I}_V)$. By assumption $U_1\times U_2$
exists and $u(U_1\times U_2) = u(U_1)\times u(U_2)$.
Consider the morphism $\phi : V \to u(U_1\times U_2)$
corresponding to $(\phi_1, \phi_2)$ by the universal property
of products. Clearly the object $\phi : V \to u(U_1\times U_2)$
maps to both $\phi_1 : V \to u(U_1)$ and $\phi_2 : V \to u(U_2)$.
\end{proof}

\noindent
Given $g : V' \to V$ in $\mathcal{D}$ we get a functor
$\overline{g} : \mathcal{I}_V \to \mathcal{I}_{V'}$
by setting $\overline{g}(U, \phi) = (U, \phi \circ g)$
on objects. Given a presheaf $\mathcal{F}$ on $\mathcal{C}$
we obtain a functor
$$
\mathcal{F}_V :
\mathcal{I}_V^{opp}
\longrightarrow
\textit{Sets}, \quad
(U, \phi)
\longmapsto
\mathcal{F}(U).
$$
In other words, $\mathcal{F}_V$ is a presheaf of sets on $\mathcal{I}_V$.
Note that we have $\mathcal{F}_{V'} \circ \overline{g} = \mathcal{F}_V$.
We define
$$
u_p\mathcal{F}(V) =
\colim_{\mathcal{I}_V^{opp}} \mathcal{F}_V
$$
As a colimit we obtain for each $(U, \phi) \in \Ob(\mathcal{I}_V)$
a canonical map $\mathcal{F}(U)\xrightarrow{c(\phi)}u_p\mathcal{F}(V)$.
For $g : V' \to V$ as above there is a
canonical restriction map
$g^* : u_p\mathcal{F}(V) \to u_p\mathcal{F}(V')$
compatible with
$\mathcal{F}_{V'} \circ \overline{g} = \mathcal{F}_V$
by Categories, Lemma \ref{categories-lemma-functorial-colimit}.
It is the unique map so that for all $(U, \phi) \in \Ob(\mathcal{I}_V)$
the diagram
$$
\xymatrix{
\mathcal{F}(U) \ar[r]^{c(\phi)} \ar[d]_{\text{id}}
&
u_p\mathcal{F}(V) \ar[d]^{g^*}
\\
\mathcal{F}(U) \ar[r]^{c(\phi \circ g)}
&
u_p\mathcal{F}(V')
}
$$
commutes. The uniqueness of these maps implies that we obtain a
presheaf. This presheaf will be denoted $u_p\mathcal{F}$.

\begin{lemma}
\label{lemma-recover}
There is a canonical map
$\mathcal{F}(U) \to u_p\mathcal{F}(u(U))$,
which is compatible with restriction maps
(on $\mathcal{F}$ and on $u_p\mathcal{F}$).
\end{lemma}

\begin{proof}
This is just the map $c(\text{id}_{u(U)})$ introduced above.
\end{proof}

\noindent
Note that any map of presheaves $\mathcal{F} \to \mathcal{F}'$
gives rise to compatible systems of maps between functors
$\mathcal{F}_V \to \mathcal{F}'_V$, and hence to a map
of presheaves $u_p\mathcal{F} \to u_p\mathcal{F}'$. In other
words, we have defined a functor
$$
u_p :
\textit{PSh}(\mathcal{C})
\longrightarrow
\textit{PSh}(\mathcal{D})
$$

\begin{lemma}
\label{lemma-adjoints-u}
The functor $u_p$ is a left adjoint to the functor $u^p$.
In other words the formula
$$
\Mor_{\textit{PSh}(\mathcal{C})}(\mathcal{F}, u^p\mathcal{G})
=
\Mor_{\textit{PSh}(\mathcal{D})}(u_p\mathcal{F}, \mathcal{G})
$$
holds bifunctorially in $\mathcal{F}$ and $\mathcal{G}$.
\end{lemma}

\begin{proof}
Let $\mathcal{G}$ be a presheaf on $\mathcal{D}$ and let
$\mathcal{F}$ be a presheaf on $\mathcal{C}$.
We will show that the displayed formula holds
by constructing maps either way. We will leave
it to the reader to verify they are each others inverse.

\medskip\noindent
Given a map $\alpha : u_p \mathcal{F} \to \mathcal{G}$
we get $u^p\alpha : u^p u_p \mathcal{F} \to u^p \mathcal{G}$.
Lemma \ref{lemma-recover} says that there is a
map $\mathcal{F} \to u^p u_p \mathcal{F}$. The composition
of the two gives the desired map. (The good thing about this construction
is that it is clearly functorial in everything in sight.)

\medskip\noindent
Conversely, given a map $\beta : \mathcal{F} \to u^p\mathcal{G}$
we get a map $u_p\beta : u_p\mathcal{F} \to u_p u^p\mathcal{G}$.
Let $V \in \Ob(\mathcal{D})$.
We claim that the functor $u^p\mathcal{G}_V$ on $\mathcal{I}_V$
has a canonical map to the constant functor with value $\mathcal{G}(V)$.
Namely, for every object $(U, \phi)$ of $\mathcal{I}_V$,
the value of $u^p\mathcal{G}_V$ on this object is $\mathcal{G}(u(U))$
which maps to $\mathcal{G}(V)$ by $\mathcal{G}(\phi) = \phi^*$.
This is a transformation of functors because $\mathcal{G}$ is a functor
itself. This leads to a map $u_p u^p \mathcal{G}(V) \to \mathcal{G}(V)$.
Another trivial verification shows that this is functorial in $V$
leading to a map of presheaves $u_p u^p \mathcal{G} \to \mathcal{G}$.
The composition $u_p\mathcal{F} \to u_p u^p\mathcal{G} \to
\mathcal{G}$ is the desired map.
\end{proof}

\begin{remark}
\label{remark-functoriality-presheaves-values}
Suppose that $\mathcal{A}$ is a category such that
any diagram $\mathcal{I}_Y \to \mathcal{A}$ has a
colimit in $\mathcal{A}$. In this case it is clear
that there are functors $u^p$ and $u_p$, defined in
exactly the same way as above, on the categories
of presheaves with values in $\mathcal{A}$.
Moreover, the adjointness of the pair
$u^p$ and $u_p$ continues to hold in this setting.
\end{remark}

\begin{lemma}
\label{lemma-pullback-representable-presheaf}
Let $u : \mathcal{C} \to \mathcal{D}$ be a functor between categories.
For any object $U$ of $\mathcal{C}$ we have $u_ph_U = h_{u(U)}$.
\end{lemma}

\begin{proof}
By adjointness of $u_p$ and $u^p$ we have
$$
\Mor_{\textit{PSh}(\mathcal{D})}(u_ph_U, \mathcal{G})
=
\Mor_{\textit{PSh}(\mathcal{C})}(h_U, u^p\mathcal{G})
=
u^p\mathcal{G}(U) =
\mathcal{G}(u(U))
$$
and hence by Yoneda's lemma we see that $u_ph_U = h_{u(U)}$ as
presheaves.
\end{proof}

















